\section{Training setup}
\label{sec:training}

\subsection{Data splits and normalization}
\label{subsec:train-splits}
Models are trained on the processed datasets detailed in Section~\ref{sec:data-generation}. To prevent scenario overlap, training, validation, and test splits are generated from independent random pools with distinct random seeds. Each solver target (hydrostatic, MUSCL-HR, and Boussinesq) comprises $10{,}000$ training, $1{,}000$ validation, and $2{,}500$ test scenarios sharing identical domain configurations to enable direct target comparisons.

Normalization statistics are computed exclusively on the training split to prevent data leakage. Bathymetry and source input channels are standardized, while trajectory targets are standardized during training and denormalized to the stored benchmark scale for metric evaluation. The \texttt{initial\_depth} feature remains in its stored scale.
\subsection{Optimization}
\label{subsec:train-optim}
Models are trained using the AdamW optimizer with an initial learning rate of $5 \times 10^{-4}$, weight decay of $10^{-6}$, and gradient clipping at a maximum norm of $1.0$. The learning rate decays via cosine annealing to a minimum of $5 \times 10^{-5}$. The batch size is set to $64$ for most direct models, but reduced to $8$ for U-FNO and $4$ for ConvLSTM due to memory constraints.

We optimize a horizon-weighted mean-squared error that places greater weight on later frames:
\begin{equation}
\mathcal{L} = \frac{\sum_{t=1}^{T} w_t \mathrm{MSE}_t}{\sum_{t=1}^{T} w_t}, \qquad w_t = w_{\min} + (w_{\max} - w_{\min})\!\left(\frac{t - 1}{T - 1}\right)^p,
\end{equation}
with $w_{\min} = 0.8$, $w_{\max} = 2.0$, and $p = 1.7$. Model selection and early stopping (patience of 20 epochs, up to a maximum of 256 epochs) are governed by the relative $L_2$ error on \emph{normalized} validation targets. Final reported metrics are subsequently evaluated after target denormalization in benchmark-scale units.


\subsection{Experiment configurations}
\label{subsec:train-configs}
The primary direct-prediction models are evaluated across Hydrostatic, MUSCL-HR, and Boussinesq targets using a unified training procedure (loss function, selection metric, and maximum epoch budget). Hydrostatic comparisons comprise FNO, F-FNO, FNO-modes8, FNO-modes20, U-FNO, WNO, CNN, U-Net, and ConvLSTM. The weakly dispersive Boussinesq model serves as a qualitative contrast under identical scenario domain settings.

For autoregressive drift analysis, Window-FNO-5 and Window-F-FNO-5 are trained on hydrostatic data to predict 5-frame blocks ($K = 5$) given two prior elevation states. In inference, rollouts are seeded with the first target frame and evaluated over the remaining 49 frames. These windowed models quantify conditional error propagation rather than direct single-pass performance.

Uncertainty quantification relies on a seven-member deep ensemble of the baseline FNO, differing strictly in initialization seed. The selected direct and windowed checkpoints use seed 18, while the ensemble uses seeds $\{11, 22, 33, 44, 55, 66, 77\}$. Sample-scaling runs also use seed 18, so their full-data comparator has the same training seed.

\subsection{Hardware and runtime setting}
\label{subsec:train-hardware}

Dataset generation and the final non-timing evaluation use university cluster resources, with CUDA model evaluation assigned to a DGX-A100 node. The study does not include a controlled training-efficiency comparison; archival wall-clock costs for the conditional windowed models are reported separately in Appendix~\ref{app:training-time} and are interpreted only as hardware- and load-dependent computational cost.

Random seed allocations are structured as follows: primary direct and windowed models use seed 18; dataset-scaling experiments use seed 18; and the seven-member uncertainty ensemble employs seeds $\{11, 22, 33, 44, 55, 66, 77\}$. Complete configuration files, training logs, and selected model states are provided for reproducibility.
